\h{Handling Errors with try and except}
Imagine your Python script is a detailed set of construction instructions for a complex project. The instructions are followed perfectly, one after the other. But what happens if the script encounters an impossible situation? What if it's told to divide a structural load by an area of zero, or to read material specifications from a file that doesn't exist?

Without a contingency plan, the entire project comes to a halt. When Python runs into an operation it cannot complete, it raises an exception—a signal that a critical error has occurred. If this exception is unhandled, it's like a "stop work" order: your entire program crashes. The try and except statements are Python's mechanism for creating a contingency plan. They allow you to anticipate potential errors and handle them gracefully, ensuring your script remains in control and provides helpful feedback instead of abruptly failing.

\begin{tcolorbox}[% passing options here
    title=Focus of this chapter:,
    colframe=orange,
    colback = white,
%    colback=orange!40!yellow!30, % think of mixing 2 colors with sliders
    arc=2mm % sharp corners
 ]
\begin{itemize}
    \item What exceptions are and where they come from
    \item The try/except pattern for handling errors
    \item Using \c{else} (happy path) and \c{finally} (always runs)
    \item Minimal, runnable examples for common beginner errors
    \item Input validation loops with friendly messages
    \item When (and how) to raise your own exceptions
\end{itemize}
\end{tcolorbox}

\hh{Basic try/except}

If a line \emph{might} fail, put it inside a \c{try} block and catch the specific error you expect. This keeps the program from crashing and gives the user a clear message.

\code{c09_basic_try}{python}{
    # Divide two numbers with protection against ZeroDivisionError
    a = 12
    b = 0

    try:
        result = a / b
        print("result:", result)
    except ZeroDivisionError:
        print("Cannot divide by zero. Please choose a non-zero denominator.")
    # Outcome:
    # Cannot divide by zero. Please choose a non-zero denominator.
}{Catch a specific exception to prevent a crash.}

\hh{Catching multiple exception types}

You can handle different problems with different messages. Use except <Type> lines in order from most specific to most general.

\code{c09_multiple_except}{python}{
    raw_width = "three point five"   # user typed a phrase, not a number
    raw_count = "0"

    try:
        width = float(raw_width)     # may raise ValueError
        count = int(raw_count)       # may raise ValueError
        avg = width / count          # may raise ZeroDivisionError
        print("average:", avg)
    except ValueError:
        print("Please enter numbers only (e.g., 3.5).")
    except ZeroDivisionError:
        print("Count cannot be zero.")
    # Outcome:
    # Please enter numbers only (e.g., 3.5).
}{Handle different errors with tailored messages.}

\hh{Access the exception object}

To show or log the original message from Python, capture the exception as e. This is useful for debugging or when you need to include details (like a filename) in your feedback.

\code{c09_exception_object}{python}{
    path = "missing_file.txt"

    try:
        with open(path, "r") as f:
            data = f.read()
    except FileNotFoundError as e:
        print("File problem:", e)
    # Outcome (message varies by OS):
    # File problem: [Errno 2] No such file or directory: 'missing_file.txt'
}{Capture and print the system message.}

\hh{else: the happy path}

Use else for code that should run only when no exception happens. It keeps the success path uncluttered and separates it from error handling.

\code{c09_else}{python}{
    raw = "24.0"

    try:
        area = float(raw)
    except ValueError:
        print("Not a number.")
    else:
        print("Area is", area, "m^2")  # runs only if no error
    # Outcome:
    # Area is 24.0 m^2
}{Place the success path in else.}

\hh{finally: always run cleanup}

Code in finally runs whether there was an error or not. Use it for cleanup actions like closing files or printing a status line.

\code{c09_finally}{python}{
    print("Starting...")

    try:
        x = 1 / 0
    except ZeroDivisionError:
        print("Handled division by zero.")
    finally:
        print("Done (this always runs).")
    # Outcome:
    # Starting...
    # Handled division by zero.
    # Done (this always runs).
}{Use finally for actions that must happen.}

\begin{NoteBox}
\textbf{Order matters.} Put specific except clauses first (e.g., \c{ValueError}), then optional broader ones (e.g., \c{Exception}). A bare \c{except:} is discouraged—use \c{except Exception as e:} if you truly need a catch-all.
\end{NoteBox}

% ====================================================================
% PART II — COMMON BEGINNER ERRORS (WITH MINIMAL EXAMPLES)
% ====================================================================

\hh{TypeError (wrong type for an operation)}

Happens when an operation receives an unexpected type (e.g., adding number to string).

\code{c09_typeerror_min}{python}{
    # This raises TypeError: cannot add float and str
    try:
        total = 3.0 + "5.0"
        print(total)
    except TypeError as e:
        print("TypeError:", e)
    # Outcome:
    # TypeError: unsupported operand type(s) for +: 'float' and 'str'
}{Mixing types can raise TypeError.}

\hh{ValueError (bad content for a type)}

Happens when converting text to a number and the text is not numeric.

\code{c09_valueerror_min}{python}{
    # This raises ValueError: 'Lobby' is not a float
    try:
        x = float("Lobby")
        print(x)
    except ValueError as e:
        print("ValueError:", e)
    # Outcome:
    # ValueError: could not convert string to float: 'Lobby'
}{Invalid numeric text causes ValueError.}

\hh{ZeroDivisionError (divide by zero)}

\code{c09_zerodiv_min}{python}{
    try:
        print(10 / 0)
    except ZeroDivisionError as e:
        print("ZeroDivisionError:", e)
    # Outcome:
    # ZeroDivisionError: division by zero
}{Guard divisions against zero.}

\hh{NameError (variable not defined)}

\code{c09_nameerror_min}{python}{
    try:
        print(level)  # 'level' was never defined
    except NameError as e:
        print("NameError:", e)
    # Outcome:
    # NameError: name 'level' is not defined
}{Use clear, consistent variable names and define them before use.}

\hh{IndexError and KeyError (missing position or key)}

\code{c09_index_key_min}{python}{
    try:
        nums = [10, 20]
        print(nums[5])              # IndexError
    except IndexError as e:
        print("IndexError:", e)

    try:
        room = {"name": "Lobby", "area": 24.0}
        print(room["level"])        # KeyError
    except KeyError as e:
        print("KeyError:", e)
    # Outcome:
    # IndexError: list index out of range
    # KeyError: 'level'
}{Check lengths and use dict.get("key") with defaults.}

\hh{FileNotFoundError (path is wrong)}

\code{c09_filenotfound_min}{python}{
    try:
        with open("no_such_file.csv", "r") as f:
            data = f.read()
    except FileNotFoundError as e:
        print("FileNotFoundError:", e)
    # Outcome:
    # FileNotFoundError: [Errno 2] No such file or directory: 'no_such_file.csv'
}{Verify paths; show friendly hints to the user.}

% ====================================================================
% PART III — FRIENDLY INPUT VALIDATION
% ====================================================================

\hh{Single conversion with message}

Do one thing: try to convert, catch failure, print a friendly message. This is the smallest building block.

\code{c09_input_once}{python}{
    raw = "3O"  # note: letter 'O', not zero
    try:
        width = float(raw)
        print("width:", width)
    except ValueError:
        print("Please enter a numeric value, e.g., 3.0")
    # Outcome:
    # Please enter a numeric value, e.g., 3.0
}{Convert with a clear error message.}

\hh{Loop until valid input (while + try/except)}

Ask repeatedly until the user gives a positive number. This is the standard pattern for beginner-friendly scripts.

\code{c09_input_loop}{python}{
    while True:
        txt = input("Width (m): ").strip()
        try:
            w = float(txt)
            if w <= 0:
                print("Please enter a number greater than zero.")
                continue
            break
        except ValueError:
            print("Please enter a numeric value, e.g., 3.2")
    print("Width accepted:", w)
}{Repeat until valid input is entered.}

\begin{NoteBox}
\textbf{Style.} Prefer “ask--try--convert” over pre-checking with \c{.isdigit()}, because \c{float()} correctly handles decimals and negatives. This “try it and handle failures” approach is called EAFP (Easier to Ask Forgiveness than Permission).
\end{NoteBox}

% ====================================================================
% PART IV — RAISING YOUR OWN EXCEPTIONS
% ====================================================================

\hh{Validate and raise}

When a function receives invalid inputs, raise a clear exception. This stops the function immediately and lets the caller decide how to respond.

\code{c09_raise_valueerror}{python}{
    def area_m2(width, length):
        if width <= 0 or length <= 0:
            raise ValueError("width and length must be positive")
        return width * length

    try:
        print(area_m2(3.0, 5.0))     # OK
        print(area_m2(-2.0, 5.0))    # raises ValueError
    except ValueError as e:
        print("Invalid input:", e)
    # Outcome:
    # 15.0
    # Invalid input: width and length must be positive
}{Use raise to signal invalid inputs.}

\hh{Wrap a call and report (caller side)}

The caller catches the error and decides what to do next (show message, ask again, etc.).

\code{c09_caller_wrap}{python}{
    def parse_positive_float(text):
        try:
            x = float(text)
        except ValueError:
            raise ValueError("Not a number")
        if x <= 0:
            raise ValueError("Must be positive")
        return x

    user_input = "zero"
    try:
        val = parse_positive_float(user_input)
        print("value:", val)
    except ValueError as e:
        print("Problem:", e)
    # Outcome:
    # Problem: Not a number
}{Functions raise; callers handle.}

% ====================================================================
% PART V — PRACTICAL FILE HANDLING
% ====================================================================

\hh{Read a file safely with messages}

Catch file errors and show a helpful hint. Use with open(...) to ensure files are closed even on errors; combine with try/except for messages.

\code{c09_file_safe}{python}{
    path = "rooms.csv"  # try a wrong path first to see the message
    try:
        with open(path, "r", encoding="utf-8") as f:
            content = f.read()
    except FileNotFoundError:
        print("Could not find", path)
        print("Tip: place the file in the same folder as your script or use an absolute path.")
    else:
        print("File read OK, length:", len(content))
}{Handle missing files and succeed gracefully.}

\hh{finally for cleanup (illustration)}

Most of the time, with handles closing for you. This example shows manual cleanup with finally for contexts that need it.

\code{c09_finally_cleanup}{python}{
    f = None
    try:
        f = open("rooms.csv", "r", encoding="utf-8")
        content = f.read()
        print("Read", len(content), "characters")
    except FileNotFoundError:
        print("rooms.csv not found")
    finally:
        if f is not None:
            f.close()  # ensures closing even if an error happened
            print("File closed (finally).")
}{Use finally when resources must be released.}

% ====================================================================
% PART VI — REFERENCE AND PRACTICE
% ====================================================================

\hh{Common exceptions (quick reference)}

\renewcommand{\arraystretch}{1.12}
\begin{table}[h]
\centering
\begin{tabularx}{\textwidth}{|l|X|X|}
\hline
\textbf{Exception} & \textbf{Typical cause} & \textbf{Beginner fix} \\ \hline
\texttt{ValueError} & Converting non-numeric text to number & Validate input; show an example format \\ \hline
\texttt{TypeError} & Mixing types (e.g., number + string) & Convert types first; keep units clear \\ \hline
\texttt{ZeroDivisionError} & Dividing by zero & Check denominator; prompt for non-zero \\ \hline
\texttt{NameError} & Variable not defined yet & Define before use; avoid typos \\ \hline
\texttt{IndexError} & List index out of range & Check length; guard with \texttt{if i < len(lst)} \\ \hline
\texttt{KeyError} & Missing dict key & Use \texttt{dict.get("key", default)} \\ \hline
\texttt{FileNotFoundError} & Wrong path or filename & Verify location; show path tips to user \\ \hline
\end{tabularx}
\caption{Frequent exceptions, causes, and quick fixes.}
\end{table}

\hh{Mini-Lab A — Robust area input}

Goal: Ask for width and length (floats $> 0$), compute area, and handle bad input without crashing.

\code{c09_minilab_area}{python}{
    def read_positive_float(prompt):
        while True:
            txt = input(prompt).strip()
            try:
                x = float(txt)
            except ValueError:
                print("Please type a number like 3.2")
                continue
            if x <= 0:
                print("Please enter a positive value")
                continue
            return x

    w = read_positive_float("Width (m): ")
    l = read_positive_float("Length (m): ")
    print("Area:", w * l, "m^2")
}{Loop plus try/except for friendly validation.}

\hh{Mini-Lab B — Safe average}

Goal: Read a list of numbers from a single line (e.g., 10 20 30), compute the average, and handle empty input and zero-length cases gracefully.

\code{c09_minilab_average}{python}{
    line = input("Enter numbers separated by spaces: ").strip()
    try:
        nums = [float(tok) for tok in line.split()]
    except ValueError:
        print("Please enter only numbers like: 10 20 30")
    else:
        if len(nums) == 0:
            print("No numbers provided.")
        else:
            print("Average:", sum(nums) / len(nums))
}{Comprehension wrapped in try/except plus checks.}

\hh{Mini-Lab C — File report}

Goal: Ask for a filename, print size (characters), and show a friendly hint if the file is missing.

\code{c09_minilab_file}{python}{
    path = input("File to open: ").strip()
    try:
        with open(path, "r", encoding="utf-8") as f:
            data = f.read()
    except FileNotFoundError:
        print("Couldn't find that file.")
        print("Tip: is it in the same folder? Did you spell the name correctly?")
    else:
        print("OK:", len(data), "characters")
}{User-friendly file handling.}

\hh{Key Takeaways}
\begin{tcolorbox}[% passing options here
    colframe=orange,
    colback = white,
%    colback=orange!40!yellow!30, % think of mixing 2 colors with sliders
    arc=2mm % sharp corners
 ]
\begin{itemize}
    \item Handle problems you expect (ValueError, ZeroDivisionError) right where they can happen.
    \item Keep success code in else, and necessary cleanup in finally.
    \item Show clear, polite messages and examples of valid input.
    \item Raise your own exceptions inside functions to enforce rules.
\end{itemize}
\end{tcolorbox}